% Bruchrechnung – bank/bruchrechnung/ (zusammenbau v0.4, Heft)
\einheitenkopf[t1]{Bruchrechnung}
\setcounter{aufgabe}{0}

% Anlauf (ohne Prüfkennung)
\begin{aufgabe}{Addieren/Subtrahieren – Anlauf}
\begin{teile}
\teil Färbe $\frac{1}{6}$, dann weitere $\frac{4}{6}$ – wie viel ist gefärbt?

\bruchrechteck{0}{6}
\leerfeld
\teil $\frac{2}{11} + \frac{5}{11}$
\teil $\frac{5}{14} + \frac{3}{14}$
\end{teile}
\end{aufgabe}

\begin{aufgabe}{Addieren/Subtrahieren – Prüfungsaufgaben}
\begin{teile}
\steil Ein Glücksrad hat 10 gleich große Felder: 5 rot, 3 blau, 2 gelb. Es wird zweimal gedreht. Die Pfade „zweimal rot“, „zweimal blau“ und „zweimal gelb“ haben die Wahrscheinlichkeiten $\frac{1}{4}$, $\frac{9}{100}$ und $\frac{4}{100}$. Wie groß ist die Wahrscheinlichkeit für zweimal dieselbe Farbe? Gib sie gekürzt an. \hfill (P10 2014 OS)
\steil In einer Tüte sind 10 Bonbons: 6 Kirsche, 3 Zitrone, 1 Apfel. Du ziehst eins, legst es zurück und ziehst noch eins. Die Pfade „zweimal Kirsche“, „zweimal Zitrone“ und „zweimal Apfel“ haben die Wahrscheinlichkeiten $\frac{9}{25}$, $\frac{9}{100}$ und $\frac{1}{100}$. Wie groß ist die Wahrscheinlichkeit für zweimal dieselbe Sorte? Gib sie gekürzt an. \hfill (P10 2014 OS)
\steil Ein Los bestimmt die Reihenfolge von 9 Kindern beim Staffellauf. Die Wahrscheinlichkeiten, dass Mia als Erste, als Zweite oder als Dritte gezogen wird, sind $\frac{1}{9}$, $\frac{8}{72}$ und $\frac{56}{504}$. Wie groß ist die Wahrscheinlichkeit, dass Mia unter den ersten drei ist? Gib sie gekürzt an. \hfill (P10 2015 OS)
\steil Beim Vorlesewettbewerb lost die Lehrerin die Reihenfolge von 12 Kindern aus. Die Wahrscheinlichkeiten, dass Jan als Erster, als Zweiter oder als Dritter liest, sind $\frac{1}{12}$, $\frac{11}{132}$ und $\frac{110}{1320}$. Wie groß ist die Wahrscheinlichkeit, dass Jan unter den ersten drei ist? Gib sie gekürzt an. \hfill (P10 2015 OS)
\steil Eine Münze wird dreimal geworfen. Die Wahrscheinlichkeit für genau zweimal Zahl ist $\frac{3}{8}$, für dreimal Zahl $\frac{1}{8}$. Wie groß ist die Wahrscheinlichkeit für mindestens zweimal Zahl? Gib sie gekürzt an. \hfill (P10 2020 OS)
\steil Sara wirft drei Freiwürfe. Die Wahrscheinlichkeit für genau zwei Treffer ist $\frac{27}{64}$, für drei Treffer ebenfalls $\frac{27}{64}$. Wie groß ist die Wahrscheinlichkeit für mindestens zwei Treffer? Gib sie gekürzt an. \hfill (P10 2020 OS)
\steil Tim hat vier mögliche Codes für sein Fahrradschloss und probiert sie nacheinander, ohne einen zu wiederholen. Die Wahrscheinlichkeit, dass der erste Versuch passt, ist $\frac{1}{4}$; dass erst der zweite passt, ist $\frac{3}{12}$. Wie groß ist die Wahrscheinlichkeit, dass das Schloss spätestens beim zweiten Versuch aufgeht? Gib sie gekürzt an. \hfill (P10 2017 OS)
\steil Nina weiß, dass ihr Passwort eines von acht möglichen ist, und probiert sie nacheinander, ohne eines zu wiederholen. Die Wahrscheinlichkeit, dass der erste Versuch passt, ist $\frac{1}{8}$; dass erst der zweite passt, ist $\frac{7}{56}$. Wie groß ist die Wahrscheinlichkeit, dass sie spätestens beim zweiten Versuch drin ist? Gib sie gekürzt an. \hfill (P10 2017 OS)
\teil Zwei Kreisel werden gedreht. Die Wahrscheinlichkeit, dass beide Rot zeigen, ist $\frac{4}{30}$. Wie groß ist die Wahrscheinlichkeit, dass nicht beide Rot zeigen? Gib sie gekürzt an. \hfill (P10 2026 FOR)
\teil Ali zieht an zwei Losbuden je ein Los. Die Wahrscheinlichkeit für zwei Gewinne ist $\frac{6}{48}$. Wie groß ist die Wahrscheinlichkeit, dass er nicht zweimal gewinnt? Gib sie gekürzt an. \hfill (P10 2026 FOR)
\teil Ein Glücksrad hat 10 gleich große Felder mit den Zahlen 1 bis 10. Die Wahrscheinlichkeit für 8, 9 oder 10 ist $\frac{3}{10}$. Wie groß ist die Wahrscheinlichkeit für eine Zahl von 1 bis 7? \hfill (P10 2014 OS)
\teil Ein Spielwürfel hat 12 Flächen mit den Zahlen 1 bis 12. Die Wahrscheinlichkeit für 1 oder 12 ist $\frac{2}{12}$. Wie groß ist die Wahrscheinlichkeit für weder 1 noch 12? Gib sie gekürzt an. \hfill (P10 2014 OS)
\end{teile}
\end{aufgabe}

% Anlauf (ohne Prüfkennung)
\begin{aufgabe}{Dezimal addieren/subtrahieren – Anlauf}
\begin{teile}
\teil Welche Zahl steht in der Stellenwerttafel?

\sachtabelle{ccc}{Einer & Zehntel & Hundertstel}{4 & 3 & 6}
\leerfeld
\teil $3{,}4 + 2{,}5$
\teil $5{,}6 + 1{,}23$
\end{teile}
\end{aufgabe}

\begin{aufgabe}{Dezimal addieren/subtrahieren – Prüfungsaufgaben}
\begin{teile}
\teil Familie Aydin hat im letzten Jahr $1\,847{,}60$ € für Heizöl bezahlt, in diesem Jahr $2\,016{,}35$ €. Zeige, dass sie in diesem Jahr $168{,}75$ € mehr bezahlt hat. \hfill (P10 2014 OS)
\teil Ein Sportverein hat im Vorjahr $912{,}40$ € für Strom bezahlt, in diesem Jahr $1\,004{,}15$ €. Zeige, dass es in diesem Jahr $91{,}75$ € mehr sind. \hfill (P10 2014 OS)
\end{teile}
\end{aufgabe}

% Anlauf (ohne Prüfkennung)
\begin{aufgabe}{Multiplizieren – Anlauf}
\begin{teile}
\teil „$\frac{1}{4}$ von 28 Kindern“ – heißt „von“ hier mal? \janein
\teil $\frac{2}{9} \cdot 4$
\teil $\frac{2}{5} \cdot \frac{3}{7}$
\end{teile}
\end{aufgabe}

\begin{aufgabe}{Multiplizieren – Prüfungsaufgaben}
\begin{teile}
\steil In einer Dose liegen 15 Murmeln: 7 grün, 5 weiß, 3 schwarz. Drei werden nacheinander ohne Zurücklegen gezogen. Tom behauptet: „Dreimal grün ist doppelt so wahrscheinlich wie dreimal weiß.“ Widerlege die Behauptung mit einer Rechnung. \hfill (P10 2019 OS)
\steil In einer Lostrommel sind 10 Lose: 6 Nieten und 4 Gewinne. Drei Lose werden nacheinander ohne Zurücklegen gezogen. Lena behauptet: „Dreimal eine Niete ist dreimal so wahrscheinlich wie dreimal ein Gewinn.“ Widerlege die Behauptung mit einer Rechnung. \hfill (P10 2019 OS)
\steil Auf einem Teller liegen 12 Muffins, 9 mit Schokolade und 3 mit Nuss. Emma nimmt nacheinander zwei, ohne hinzusehen. Wie groß ist die Wahrscheinlichkeit, dass beide mit Schokolade sind? Gib sie gekürzt an. \hfill (P10 2018 OS)
\steil In einer Dose sind 10 Kekse, 7 Butterkekse und 3 Schokokekse. Ole nimmt nacheinander zwei, ohne hinzusehen. Wie groß ist die Wahrscheinlichkeit, dass beide Butterkekse sind? Gib sie gekürzt an. \hfill (P10 2018 OS)
\steil In einem Beutel sind 14 Kugeln, 4 davon gelb. Drei Kugeln werden ohne Zurücklegen gezogen. Der Pfad gelb–gelb–gelb hat die Wahrscheinlichkeit $\frac{4}{14} \cdot \frac{3}{13} \cdot \frac{2}{12}$. Berechne sie gekürzt und entscheide, ob sie kleiner als $1\,\%$ ist. \hfill (P10 2019 OS)
\steil In einer Lostrommel sind 16 Lose, 3 davon Gewinne. Drei Lose werden ohne Zurücklegen gezogen. Der Pfad Gewinn–Gewinn–Gewinn hat die Wahrscheinlichkeit $\frac{3}{16} \cdot \frac{2}{15} \cdot \frac{1}{14}$. Berechne sie gekürzt und entscheide, ob sie kleiner als $1\,\%$ ist. \hfill (P10 2019 OS)
\steil Sechs Karten liegen verdeckt, eine davon ist die Gewinnkarte. Anna zieht zuerst eine Karte, danach Ben eine der übrigen fünf. Anna zieht keine Gewinnkarte mit der Wahrscheinlichkeit $\frac{5}{6}$, danach zieht Ben sie mit $\frac{1}{5}$. Wie groß ist die Wahrscheinlichkeit, dass Ben die Gewinnkarte zieht? \hfill (P10 2024 OS)
\steil In einer Schachtel sind 10 Lose, 2 davon Freikarten. Anna zieht zuerst, dann Ben. Anna zieht keine Freikarte mit der Wahrscheinlichkeit $\frac{8}{10}$, danach zieht Ben eine mit $\frac{2}{9}$. Wie groß ist die Wahrscheinlichkeit, dass Anna keine und Ben eine Freikarte zieht? Gib sie gekürzt an. \hfill (P10 2024 OS)
\end{teile}
\end{aufgabe}

% Anlauf (ohne Prüfkennung)
\begin{aufgabe}{Dezimal multiplizieren/dividieren – Anlauf}
\begin{teile}
\teil $3{,}72 \cdot 10$
\teil $45{,}7 : 10$
\end{teile}
\end{aufgabe}

\begin{aufgabe}{Dezimal multiplizieren/dividieren – Prüfungsaufgaben}
\begin{teile}
\teil Ein Lieferwagen braucht 8 l Diesel auf 100 km und fährt im Jahr $25\,000$ km. Ein Liter Diesel kostet $162{,}4$ ct. Berechne die Kraftstoffkosten für ein Jahr in Euro. \hfill (P10 2014 OS)
\teil Ein Mofa braucht 3 l Benzin auf 100 km und fährt im Jahr $4\,000$ km. Ein Liter Benzin kostet $178{,}5$ ct. Berechne die Benzinkosten für ein Jahr in Euro. \hfill (P10 2014 OS)
\teil Ein Handy einmal voll zu laden kostet $0{,}35$ ct Strom. Wie viel kostet es, das Handy 30-mal zu laden? Gib das Ergebnis in Cent an. \hfill (P10 2024 OS)
\teil Eine LED-Lampe verbraucht in einer Stunde Strom für $0{,}26$ ct. Sie brennt 24 Stunden. Wie viel kostet das? Gib das Ergebnis in Cent an. \hfill (P10 2024 OS)
\teil Welcher Wert ist am kleinsten: $2{,}2$; $0{,}22$; $0{,}2^2$; $22\,\%$? \hfill (P10 2023 OS) \\ \leerfeld
\teil Welcher Wert ist am kleinsten: $5{,}5$; $0{,}55$; $0{,}5^2$; $55\,\%$? \hfill (P10 2023 OS) \\ \leerfeld
\end{teile}
\end{aufgabe}

% Anlauf (ohne Prüfkennung)
\begin{aufgabe}{Punkt vor Strich – Anlauf}
\begin{teile}
\teil $7 + 3 \cdot 5$ – kreise die Rechnung ein, die zuerst drankommt.
\teil $5 + 4 \cdot 6$
\teil $(3{,}5 + 1{,}5) \cdot 4$
\end{teile}
\end{aufgabe}

\begin{aufgabe}{Punkt vor Strich – Prüfungsaufgaben}
\begin{teile}
\teil Berechne den Wert des Terms $\frac{a + b}{c}$ für $a = 4{,}5$, $b = 1{,}5$ und $c = 4$. \hfill (P10 2021 OS)
\teil Berechne den Wert des Terms $\frac{a + b}{c}$ für $a = \frac{2}{3}$, $b = \frac{1}{6}$ und $c = 5$. \hfill (P10 2021 OS)
\teil Berechne den Wert des Terms $(a + b) : c$ für $a = 7$, $b = 2{,}6$ und $c = 1{,}2$. \hfill (P10 2016 OS)
\teil Berechne den Wert des Terms $(a + b) : c$ für $a = 10$, $b = 3{,}5$ und $c = 0{,}5$. \hfill (P10 2016 OS)
\teil Im Schwimmbad kostet der Eintritt für Erwachsene 8,00 €, für Kinder 4,00 €. Montags gibt es $25\,\%$ Rabatt. Familie Berg kommt mit Mutter, Vater, Tochter (9 Jahre) und Oma. Gibt der Term den Rabatt richtig an? \hfill (P10 2015 OS)\\
$\frac{1}{4} \cdot (24\text{ €} + 4\text{ €})$ \janein\\
$(2 \cdot 8\text{ €} + 2 \cdot 4\text{ €}) \cdot \frac{25}{100}$ \janein\\
$25 \cdot (8\text{ €} + 8\text{ €} + 8\text{ €} + 4\text{ €}) : 100$ \janein
\teil Im Kino kostet eine Karte für Erwachsene 9,50 €, für Kinder 6,00 €. Mittwochs gibt es $20\,\%$ Rabatt. Familie Kaya kommt mit Vater, Opa und zwei Kindern (10 und 12 Jahre). Gibt der Term den Rabatt richtig an? \hfill (P10 2015 OS)\\
$(9{,}50\text{ €} + 3 \cdot 6\text{ €}) \cdot 0{,}2$ \janein\\
$\frac{1}{5} \cdot (19\text{ €} + 12\text{ €})$ \janein\\
$20 \cdot (9{,}50\text{ €} + 9{,}50\text{ €} + 6\text{ €} + 6\text{ €}) : 100$ \janein
\end{teile}
\end{aufgabe}
